\begin{tikzpicture}[
  x=1cm,
  y=1cm,
  route/.style={font=\sffamily\bfseries\large, align=center},
  note/.style={font=\sffamily\small, align=center, text=black!70},
  molname/.style={font=\sffamily\bfseries\small, align=center},
  molnote/.style={font=\sffamily\scriptsize, align=center, text=black!70},
  arrow/.style={->, >=stealth, line width=0.7pt, black!85}
]

% Image sizes.  These are intentionally defined here so the layout can be
% adjusted directly from LaTeX.
\def\parentw{0.19\linewidth}
\def\parentwide{0.17\linewidth}
\def\molw{0.15\linewidth}
\def\molwb{0.15\linewidth}

% Molecule images.  Draw these first because the PNGs have white backgrounds.
\node[inner sep=0pt] (pt10parent) at (-3.30,2.75)
  {\includegraphics[width=\parentw]{pt10.png}};
\node[inner sep=0pt] (pt5ge5parent) at (3.45,2.75)
  {\includegraphics[width=\parentwide]{pt5ge5_orig.png}};

\node[inner sep=0pt] (pt6ge4) at (-5.15,-0.10)
  {\includegraphics[width=\molw]{pt6ge4.png}};
\node[inner sep=0pt] (pt5ge5) at (-3.25,-0.10)
  {\includegraphics[width=\molw]{Pt5Ge5_fromPt10Parent.png}};
\node[inner sep=0pt] (pt4ge6) at (-1.15,-0.10)
  {\includegraphics[width=\molw]{pt4ge6.png}};
\node[inner sep=0pt] (pt10frompt5ge5) at (3.45,-0.10)
  {\includegraphics[width=\molwb]{pt10_from_pt5ge5.png}};

% Route arrows.
\draw[arrow] (-3.82,1.18) -- (-4.90,0.80);
\draw[arrow] (-3.30,1.18) -- (-3.25,0.82);
\draw[arrow] (-2.78,1.18) -- (-1.55,0.80);
\draw[arrow] (3.45,1.18) -- (3.45,0.85);

% Headers and labels.  Draw these last so they remain visible over the PNG
% image boxes.
\node[route] at (-3.30,4.58) {\ce{Pt10} parent framework};
\node[note] at (-3.30,4.16) {Ge substitution followed by optimization};
\node[route] at (3.45,4.58) {\ce{Pt5Ge5} parent framework};
\node[note] at (3.45,4.16) {Ge-to-Pt substitution followed by optimization};

\node[molname] at (-3.30,1.62) {\ce{Pt10} parent};
\node[molnote] at (-3.30,1.29) {reported framework};
\node[molname] at (3.45,1.62) {\ce{Pt5Ge5} parent};
\node[molnote] at (3.45,1.29) {reported framework};

\node[molname] at (-5.15,-1.18) {\ce{Pt6Ge4}};
\node[molnote] at (-5.15,-1.51) {from \ce{Pt10}};
\node[molname] at (-3.25,-1.18) {\ce{Pt5Ge5}};
\node[molnote] at (-3.25,-1.51) {from \ce{Pt10}};
\node[molname] at (-1.35,-1.18) {\ce{Pt4Ge6}};
\node[molnote] at (-1.35,-1.51) {from \ce{Pt10}};
\node[molname] at (3.45,-1.18) {\ce{Pt10}};
\node[molnote] at (3.45,-1.51) {from \ce{Pt5Ge5}};

\node[note] at (-0.70,-2.22)
  {Upper row: reported parent frameworks. Lower row: optimized crossed-parent derivatives.};

\end{tikzpicture}
